% Conversão fiel de chapter-05-text-h-0075; origem RELATORIO_ATUAL.odt.
\hypertarget{__RefHeading___Toc4799_3118176633}{}
\refstepcounter{section}
\section*{5.1 Introdução}
\label{ages-iii:text-h-0075}
\addcontentsline{toc}{section}{5.1 Introdução}

% SOURCE text-p-1022
\originaltab{}A tecnologia RFID é amplamente utilizada em ambientes que necessitam de controle de acesso rápido e seguro, como estacionamentos corporativos\space\cite{ages-iii-rfid-parking-use-cases}, condomínios residenciais e empresas. Nesses contextos, seu uso normalmente está concentrado na identificação de veículos autorizados e na liberação de entradas e saídas. No entanto, os dados gerados por esse processo também podem ser utilizados de forma estratégica, permitindo uma gestão mais eficiente da operação.\par

% SOURCE text-p-1023
Dentro dessa proposta, o presente projeto tem como objetivo integrar um leitor RFID a uma plataforma web de gerenciamento, capaz de centralizar informações e oferecer recursos de monitoramento inteligente. A solução contempla um \textit{dashboard} administrativo com métricas de fluxo de veículos, ocupação de vagas e histórico de acessos, além de permitir interações diretas com o equipamento, como ajustes de configuração e administração de seu funcionamento.\par

% SOURCE text-p-1024
Outro diferencial do projeto está em sua natureza técnica, por envolver não apenas desenvolvimento de software, mas também integração com hardware real, o Impinj R700\space\cite{ages-iii-impinj-r700}. Dentro do contexto da AGES, trata-se de uma experiência pouco comum, exigindo do time conhecimentos de arquitetura, comunicação entre sistemas e operação de dispositivos físicos. Dessa forma, o projeto proporciona uma vivência próxima de desafios encontrados no mercado, onde soluções completas dependem da integração entre software e hardware.\par

% FIGURE draw-image-0016
\memorialfigure{imagens/original/asset-0015.jpg}{\originaltab{}\textbf{Figura 16 –} Foto do Time - AGES III}{ages-iii:draw-image-0016}{Fonte: pendente de identificação}{0.92} % <<ASSET draw-image-0016 asset-0015>>
